\section{The actual third transition and its revived second parent}
\label{tt:section}

The released transition is the full physical interval
$[t_3^a,t_3^a+\Gamma_3^{-1}]$, with
$\dot\alpha=(1-\eta_b)\mathcal F_2+\eta_b\mathcal F_3$ and
$\eta_b=\eta(\Gamma_3(t-t_3^a))$; see source (3.4)--(3.6),
Section 3.7.5, pp.~30--31, and (5.22), p.~46.
The following calculation retains that control, both evolving parents,
their physical diffusion, and the original laboratory gravity.
It does not use the exceptional fixed geometry of transition two.
All older activation weights are one on this interval.

\subsection{Exact determinant coordinates, coefficients, and inverse}

Keep $A_0,\mu,\nu=\sqrt{A_0}$, $K_j=\mu\widehat\lambda_j$,
$s_*$, and the actual inherited scales $\sigma_1,\sigma_2$.
Physical viscosity and thermal diffusivity equal $d>0$.
Write
\begin{equation}\label{tt:constants}
 a=\sin s_2,\quad c=A_0\cos s_*,\quad C=A_0\sin s_*,
 \quad s=s_3=L_3\sigma_1/\sigma_2,\quad
 \Gamma=\Gamma_3=\sigma_2\sin s.
\end{equation}
Thus $a,s$ and $d$ are distinct parameters. Let $t_b=t_2^b$,
let $r_b=|\zeta_2(t_b)|$, and put $b=r_b\sin s>0$.
The original third datum is $\zeta_3(t_b)=e(s)$ and
$M_3(t_b)=I$, while $\zeta_2(t_b)=r_be_2$.
For $e(\phi)=(\sin\phi,\cos\phi)$ and
$J=\left(\begin{smallmatrix}0&-1\\1&0\end{smallmatrix}\right)$,
the two adjacent oriented determinants are exactly
\begin{equation}\label{tt:determinants}
 \det(\zeta_1,\zeta_2)=-a,\qquad
 \det(\zeta_2,\zeta_3)=-b.
\end{equation}
Both identities persist through growth, transition and later controls
whenever their equations exist. Indeed $|\zeta_1|=1$ and
$\zeta_1$ also solves the transport equation of $\zeta_2$:
the added first shear has transpose annihilating $\zeta_1$.
Likewise the added second shear has transpose annihilating $\zeta_2$,
so $\zeta_2,\zeta_3$ solve the same third transport equation.
The determinant of two solutions of $\dot z=-D^Tz$ has derivative
$-\operatorname{tr}(D)$ times itself, by differentiating its two
columns. Every parent matrix has zero trace, proving both assertions.

Use the following coordinates, with their original dimensional factors:
\begin{equation}\label{tt:coordinates}
\begin{gathered}
 \zeta_2=h\zeta_1-aJ\zeta_1,\qquad
 \zeta_3=x\zeta_1-yJ\zeta_1,\qquad D=h^2+a^2,\\
 g=\zeta_2\cdot\zeta_3=hx+ay,\qquad
 hy-ax=b,\qquad
 x=\frac{hg-ab}{D},\quad y=\frac{ag+hb}{D},\\
 E_1=-K_1\Theta_1=A_1,\qquad
 E_2=-K_2\Theta_2=A_2/\sqrt D,\qquad
 R=|\zeta_3|^2=\frac{g^2+b^2}{D}.
\end{gathered}
\end{equation}
Here $D$ and $R$ are scalars, not velocity-gradient matrices.
In particular the actual second amplitude is $A_2=\sqrt D E_2$;
replacing it by $E_2$ would lose the evolving phase length.
Put $\theta=\phi_2$, so the full phase and angular inverse is
\begin{align}\label{tt:inverse}
 \zeta_1&=\frac1{\sqrt D}
    (h\sin\theta-a\cos\theta,\ h\cos\theta+a\sin\theta),\\
 \zeta_2&=\sqrt D(\sin\theta,\cos\theta),\nonumber\\
 \zeta_3&=\frac1{\sqrt D}
    (g\sin\theta+b\cos\theta,\ g\cos\theta-b\sin\theta),\nonumber\\
 \phi_1&=\theta-\arctan(a/h),\qquad
 \alpha=s_*-\phi_1.\nonumber
\end{align}
The domain used below has $h,g>0$ and
$g\cos\theta-b\sin\theta>0$; hence these expressions select the
positive second-component branches. The original deformation matrices are
\begin{equation}\label{tt:matrices}
 M_1=R_\alpha,\qquad
 M_2=R_{\alpha-s_*}
 \begin{pmatrix}1&-(h-\cos s_2)/a\\0&1\end{pmatrix},\qquad
 M_3=\mathcal L(t)^{-T}\mathcal L(t_b)^T,
 \quad\mathcal L=(\zeta_2\ \zeta_3).
\end{equation}
The first two preserve their original activation data, and
$\det\mathcal L=-b\ne0$. Differentiating the last formula using
$\dot{\mathcal L}=-D_{<3}^T\mathcal L$ gives
$\dot M_3=D_{<3}M_3$, $M_3(t_b)=I$, and
$M_3^{-T}e(s)=\zeta_3$. All three have determinant one.
Thus the inverse gives the complete original matrix evolution, rather
than merely the newest phase direction.

Define the exact temperature numerators
\begin{equation}\label{tt:numerators}
 N_2=a(E_1+c)+Ch,\qquad
 N_3=(E_1+c)y+Cx+bE_2.
\end{equation}
The full older physical equations are
\begin{align}\label{tt:physical_closed}
 \dot h&=a\Omega_1,&
 \dot g&=\Omega_1\frac{2ahg+(h^2-a^2)b}{D}+b\Omega_2,\\
 \dot E_1&=-C\Omega_1-dK_1^2E_1,&
 \dot\Omega_1&=-E_1\zeta_{1,1}-dK_1^2\Omega_1,\nonumber\\
 \dot E_2&=-\frac{N_2}{D}\Omega_2-dK_2^2DE_2,&
 \dot\Omega_2&=-E_2\sqrt D\sin\theta-dK_2^2D\Omega_2,\nonumber\\
 \dot\theta&=-\dot\alpha-\frac{a^2\Omega_1}{D}.&&\nonumber
\end{align}
In terms of the actual amplitude $A_2$ the fifth equation is exactly
\[
 \dot A_2=\left(\frac{ah\Omega_1}{D}-dK_2^2D\right)A_2
               -\frac{N_2}{\sqrt D}\Omega_2.
\]
The exact source feedbacks and their factored difference are
\begin{align}\label{tt:feedback}
 \mathcal F_2&=\frac{a\Omega_1(h\tan\theta-a)}{D},\\
 \mathcal F_3&=
 \frac{\Omega_1y(h\sin\theta-a\cos\theta)+b\Omega_2\sin\theta}
      {g\cos\theta-b\sin\theta},\nonumber\\
 \mathcal F_3-\mathcal F_2
 &=\frac{b\Omega_1(h\sin\theta-a\cos\theta)
                  (h\cos\theta+a\sin\theta)}
         {D\cos\theta(g\cos\theta-b\sin\theta)}
     +\frac{b\Omega_2\sin\theta}{g\cos\theta-b\sin\theta}.\nonumber
\end{align}
With $\eta_b$ as specified, the final equation in
\eqref{tt:physical_closed} becomes
\begin{equation}\label{tt:angle}
 \dot\theta=-\frac{ah\Omega_1}{D}\tan\theta
                  -\eta_b(\mathcal F_3-\mathcal F_2).
\end{equation}
At the actual transition entry $\theta=\Omega_2=0$,
$\mathcal F_3-\mathcal F_2=-abh\Omega_1/(Dg)<0$.
The switching therefore drives a positive second angle and a negative
second vorticity, as proved on the full interval below.

\begin{proof}[Verification of the closed equations and exact invariants]
The rotating base and the sum of inherited gradients satisfy
\begin{gather*}
 R_\alpha e_2=\cos s_*\zeta_1+\sin s_*J\zeta_1,\\
 G_{<3}=-(E_1+c+hE_2)\zeta_1-(C-aE_2)J\zeta_1.
\end{gather*}
Taking $-J\zeta_2\cdot G_{<2}$ and
$-J\zeta_3\cdot G_{<3}$ gives respectively $N_2,N_3$;
the shear-amplitude coefficients are therefore
$a_2=N_2/(K_2D)$ and $a_3=N_3/(K_3R)$, with their original
signs. Since $\dot\zeta_1=\dot\alpha J\zeta_1$, the phase
equations in this moving orthonormal basis give, component by component,
\[
 \dot h=a\Omega_1,\qquad
 \dot x=\Omega_1y+\frac{bh}{D}\Omega_2,\qquad
 \dot y=\frac{ba}{D}\Omega_2.
\]
The derivative of $hy-ax$ is zero, and differentiating $g=hx+ay$
proves its displayed row. The damped amplitude rows then give every
remaining row of \eqref{tt:physical_closed}. Substituting these phases
in the source feedback formula proves \eqref{tt:feedback}; the polar
phase equation gives \eqref{tt:angle}. Conversely these calculations
reverse: differentiating \eqref{tt:inverse}, with
\eqref{tt:physical_closed} and \eqref{tt:angle}, gives the three
physical phase equations and the source angular control. Together with
\eqref{tt:matrices}, the temperature inverse and the vorticity rows,
this proves equivalence on the stated domain.

There is an exact cancellation despite the revival of $\Omega_2$:
\begin{equation}\label{tt:numerator_derivatives}
 \dot N_2=-adK_1^2E_1,\qquad
 \dot N_3=-d(K_1^2E_1y+bK_2^2DE_2).
\end{equation}
Indeed
\[
\begin{aligned}
 \dot N_3={}&(-C\Omega_1-dK_1^2E_1)y
 +(E_1+c)\frac{ba}{D}\Omega_2
 +C\left(\Omega_1y+\frac{bh}{D}\Omega_2\right)\\
 &+b\left(-\frac{a(E_1+c)+Ch}{D}\Omega_2-dK_2^2DE_2\right),
\end{aligned}
\]
whose two shear groups cancel exactly. Differentiating $N_2$ gives
its first formula. Thus both numerators are conserved when $d=0$
even though the individual phases, amplitudes and controls evolve.
For positive $d$ the two losses in $\dot N_3$ are retained explicitly.
\end{proof}

\subsection{The general aggregate identity, including activation}

The cancellation has an exact form for any finite active family.
Write $B_j=w_jK_j\Theta_j\zeta_j$, $r_j=|\zeta_j|$,
$E_j=-K_j\Theta_j$, and
\[
 G_{<q}=-A_0R_\alpha e_2+\sum_{j<q}B_j,\qquad
 D_{<q}=\dot\alpha J+
      \sum_{j<q}w_j\Omega_j\frac{J\zeta_j\otimes\zeta_j}{r_j^2}.
\]
Use the phase equations and the physical temperature rows with
thermal diffusivity $d$, without added amplitude forcing. Direct
pairwise cancellation gives
\begin{equation}\label{tt:aggregate}
 \dot G_{<q}+D_{<q}^TG_{<q}
 =\sum_{j<q}(\dot w_jK_j\Theta_j
                  -dw_jK_j^3r_j^2\Theta_j)\zeta_j.
\end{equation}
For precision, the $j$th shear applied to all lower-index gradient
terms is
$w_j\Omega_j\zeta_j(J\zeta_j\cdot G_{<j})/r_j^2$.
It cancels $w_jK_j\dot\Theta_j\zeta_j$ apart from diffusion.
The terms with lower shears cancel the phase derivative in $B_j$,
its self-shear vanishes because $J\zeta_j\cdot\zeta_j=0$,
and the rotational base and phase terms cancel exactly.
The only terms left are precisely those displayed in
\eqref{tt:aggregate}, including $\dot w_j$.

Put $N_q=-J\zeta_q\cdot G_{<q}$ and
$\Delta_{jq}=\det(\zeta_j,\zeta_q)$.
For any trace-free $2\times2$ matrix $D$, direct multiplication gives
$JD^T+DJ=0$. Differentiating the bilinear pairing using
$\dot\zeta_q=-D_{<q}^T\zeta_q$ thus gives
\begin{equation}\label{tt:general_numerator}
 \dot N_q=
 \sum_{j<q}(\dot w_j-dw_jK_j^2r_j^2)
               E_j(-\Delta_{jq}).
\end{equation}
Equivalently its diffusion part is
$\sum_{j<q}dK_j^2r_j^2J\zeta_q\cdot B_j$.
The sign is fixed by $\det(\zeta_q,\zeta_j)=-\Delta_{jq}$,
not by a convention about unoriented angle sizes.
No nonadjacent determinant has been assumed constant.
For $q=3$, $-\Delta_{13}=y$, $-\Delta_{23}=b$, and $w_1=w_2=1$;
\eqref{tt:general_numerator} gives exactly
\eqref{tt:numerator_derivatives}. During activation the term
$\dot w_j E_j(-\Delta_{jq})$ remains, so inviscid constancy is
asserted only after those weights become constant.

\subsection{Full transition from the actual third-growth endpoint}

Choose the original source data and positive physical diffusivity by
\eqref{tge:Y3}--\eqref{tge:d3}, and construct the actual third threshold
in Theorem~\ref{tge:actual_threshold}.
Here are the actual endpoint bounds used in the finite proof. They are
quantities at the already constructed threshold $t_a=t_3^a$, not
assumptions about the subsequent unknown trajectory. The preceding
third-growth calculation starts at the proved second return and gives
\begin{equation}\label{tt:entrybox}
\begin{gathered}
 1\leq r_b\leq\sqrt{10},\quad
 h(t_a)\in[\sqrt{1-a^2},3),\quad 1\leq D(t_a)<10,\quad
 \theta(t_a)=\Omega_2(t_a)=0,\\
 E_1(t_a)/\sigma_1^2\in[1/4,2],\quad
 0<\Omega_1(t_a)/\sigma_1<9,\quad
 A_2(t_a)/\sigma_2^2\in[2/5,2],\\
 1\leq R(t_a)\leq2,\quad
 \tfrac12\sin s\leq\zeta_{3,1}(t_a)\leq\sin s,\quad
 \zeta_{3,2}(t_a)\geq\cos s>15/16,\\
 \frac25\leq\frac{N_3(t_a)}{\sigma_2^2\sin s}\leq\frac{16}{15},
 \qquad \frac12\leq
 \frac{\sigma_2\Omega_3(t_a)}{K_3\Theta_3(t_a)}\leq\sqrt5,\\
 \Theta_3(t_a)=-\frac{A_0}{\mu}\widehat\lambda_3^{-7/8},\qquad
 0<c/\sigma_1^2\leq1/4,\quad
 0<C/\sigma_1^2\leq a/32.
\end{gathered}
\end{equation}
In particular $E_2(t_a)/\sigma_2^2\in[1/8,2]$:
the lower bound follows from $(2/5)/\sqrt{10}>1/8$.
Also $g(t_a)=\sqrt{D(t_a)}\zeta_{3,2}(t_a)$ lies in
$[15/16,\sqrt{20}]$. These are the exact converted endpoint
bounds needed below. Retain the explicit source restrictions
\begin{equation}\label{tt:source_box}
 0<a\leq1/512,\quad 0<s\leq a,\quad L_3\geq16384,
 \qquad \delta_2=\frac{dK_2^2}{\Gamma}\leq\frac1{128},
 \qquad \delta_3=\frac{dK_3^2}{\Gamma}\leq\frac1{108}.
\end{equation}
Every bound in \eqref{tt:entrybox} and \eqref{tt:source_box} follows
from that explicit construction, as we now verify. The second hold
gives $E_1/\sigma_1^2>191/512>1/4$ and all its other displayed
first-parent bounds. Its second-temperature diffusion exponent on
the third growth is at most $\log2$, while its length increases.
Consequently
\[
 \frac{A_2(t_a)}{\sigma_2^2}
 =\frac{A_{2,b}}{\sigma_2^2}
    \frac{\sqrt{D(t_a)}}{r_b}
    \exp\!\left(-dK_2^2\int_{t_b}^{t_a}D\,dt\right)
 \in[2/5,2].
\]
Here the lower bound uses $A_{2,b}/\sigma_2^2\geq4/5$,
$\sqrt D/r_b\geq1$ and $e^{-\log2}=1/2$; the upper bound uses
$A_{2,b}/\sigma_2^2\leq1$ and $\sqrt D/r_b<2$.
Equations~\eqref{tge:horizon_geometry}, \eqref{tge:coeffbox} and
\eqref{tge:initial_rate} give the displayed third length, first
component and numerator bounds. The positive second component has
square $R-\zeta_{3,1}^2\geq1-\sin^2s=\cos^2s$, so it is at
least $\cos s$. Equation~\eqref{tge:quotient} supplies the actual
signed pair and its quotient. The two base bounds are exactly
\eqref{ave:C_bounds} after the stated change of notation.
Finally \eqref{tge:Y3} gives $L_3>16384$ and $a<1/12288<1/512$,
and its proved angle satisfies $s\leq a$. The two positive entries
in \eqref{tge:d3}, with \eqref{tge:rate_scale}, give
\[
 \delta_2\leq\frac{\log2}{150L_3}<\frac1{128},\qquad
 \delta_3\leq\frac1{120}<\frac1{108}.
\]
Thus no assumption on an unconstructed transition is present.

\begin{theorem}\label{tt:transition}
For every original source datum, with frequencies and positive physical
diffusivity chosen by \eqref{tge:Y3}--\eqref{tge:d3}, the actual third
threshold constructed above is followed by the entire original third
transition, which exists
through $t_1=t_a+\Gamma^{-1}$. Put
$\tau=\Gamma(t-t_a)$, $u=\theta/s$, and $V=\Omega_2/\sigma_2$.
Throughout $0\leq\tau\leq1$,
\begin{equation}\label{tt:revivalbounds}
 0\leq u\leq\frac{640a}{L_3},\qquad
 -\frac{7040a}{L_3}\leq V\leq0.
\end{equation}
Both inequalities have strict signs after the switching profile has
become positive on an interval. The full older geometry obeys
\begin{equation}\label{tt:olderbounds}
\begin{gathered}
 \sqrt{1-a^2}\leq h<13/4,\quad 1\leq D<11,\quad
 3/4<g<5,\quad x,y>0,\\
 1/8<E_1/\sigma_1^2\leq2,\quad
 0<\Omega_1/\sigma_1<10,\quad
 1/16<E_2/\sigma_2^2<3,\\
 \cos\theta>15/16,\quad
 g\cos\theta-b\sin\theta>2/3.
\end{gathered}
\end{equation}
All source feedback denominators are therefore positive, and
\eqref{tt:matrices} gives finite original matrices and inverses.
The third numerator and phase length satisfy
\begin{equation}\label{tt:coefficientvariation}
 e^{-11\delta_2\tau}N_3(t_a)\leq N_3(t)\leq N_3(t_a),
 \qquad
 \left|\log\frac{|\zeta_3(t)|}{|\zeta_3(t_a)|}\right|
       \leq\frac{46000a}{L_3}<\frac1{128}.
\end{equation}
The actual newest amplitude remains negative and its quotient
$v=\sigma_2\Omega_3/(K_3\Theta_3)$ satisfies $1/2\leq v\leq3$.
Its full transition logarithmic gain is bounded by
\begin{equation}\label{tt:gain}
 \frac1{18}\leq\log\frac{-\Theta_3(t_1)}{-\Theta_3(t_a)}\leq4.
\end{equation}
The exact exponent, including physical diffusion, is given in the proof.
This theorem completes the transition; it does not assert a third
shooting root or an infinite fixed-viscosity cascade.
\end{theorem}

\begin{proof}
Set $L=L_3$, $\epsilon=\sigma_1/\Gamma=s/(L\sin s)\leq2/L$,
$q=\sigma_1/\sigma_2=s/L$, $E=E_1/\sigma_1^2$,
$F=E_2/\sigma_2^2$, $\omega=\Omega_1/\sigma_1$,
$\bar c=c/\sigma_1^2$, $\bar C=C/\sigma_1^2$ and
$\delta_1=dK_1^2/\Gamma\leq\delta_2$.
Primes now mean $d/d\tau$ exclusively. The exact dimensionless
older equations, retaining each physical ratio, are
\begin{align}\label{tt:scaled}
 h'&=a\epsilon\omega,&
 g'&=\epsilon\omega\frac{2ahg+(h^2-a^2)b}{D}
                    +\frac b{\sin s}V,\\
 E'&=-\epsilon\bar C\omega-\delta_1E,&
 \omega'&=-\epsilon E\frac{h\sin(su)-a\cos(su)}{\sqrt D}
                      -\delta_1\omega,\nonumber\\
 F'&=-q\epsilon\frac{a(E+\bar c)+\bar Ch}{D}V
                      -\delta_2DF,&
 V'&=-F\sqrt D\frac{\sin(su)}{\sin s}-\delta_2DV.\nonumber
\end{align}
The exact $u'$ is \eqref{tt:angle} divided by $\Gamma s$,
with \eqref{tt:feedback}; no approximation of either feedback occurs.

Initially impose the temporary closed bounds
$0\leq u\leq1/1024$, $-1/32\leq V\leq0$,
$\sqrt{1-a^2}\leq h\leq13/4$, $3/4\leq g\leq5$,
$1/8\leq E\leq3$, $0\leq\omega\leq10$, and $1/16\leq F\leq3$.
The sign boundaries are invariant: at $u=0$, its vector field equals
$\eta b\epsilon\omega ah/(sDg)\geq0$; at $V=0$ its vector field
is nonpositive. Also
$h\sin(su)<a\cos(su)$, since $s\leq a$ and $u\leq1/1024$.
The $\omega$ equation and its positive datum imply $\omega>0$.
These conclusions can equally be obtained from their integrating
factors and a first-exit argument, so tangency of a sign boundary
causes no loss of the sign assertion.

The temporary bounds give $D\geq1$, because $h'\geq0$ and
$D(t_a)\geq1$, and $D<11$. They give
$\cos(su)>15/16$ and
\[
 g\cos(su)-b\sin(su)
 >\frac34\frac{15}{16}-\frac{4a^2}{1024}>\frac23,
\]
using $b\leq\sqrt{10}\sin s<4s\leq4a$.
Thus every equation is smooth on a fixed neighborhood of this box.
Because $E$ decreases and $\omega'\leq\epsilon E a\leq6a/L$,
\[
 E\leq2,\quad 0<\omega\leq9+6a/L<10,\quad
 h\leq3+20a/L<13/4.
\]
Integrating the $E$ row, using the temporary upper bounds, gives
\[
 E\geq\frac14-\frac{20a}{32L}-\frac3{128}>\frac18.
\]
The numerator in the $F$ equation obeys
$0<a(E+\bar c)+\bar Ch<4a$. Since $V\leq0$,
$F'\geq-11\delta_2F$, hence
$F\geq\tfrac18e^{-11/128}>1/16$.
Its positive derivative contribution is at most
$q\epsilon(4a)/32\leq as/(4L^2)$, so $F<3$.
The older contribution to $g'$ is positive, and
\[
 g\geq\frac{15}{16}-\frac{\sqrt{10}}{32}>\frac34,
 \qquad
 g\leq\sqrt{20}+\frac{300a}{L}<5.
\]
For the upper estimate use
$[2ahg+(h^2-a^2)b]/D\leq2ag/h+b<15a$
and $\epsilon\omega\leq20/L$.
Also $hg-ab>0$ and $ag+hb>0$, proving $x,y>0$.

It remains to strictly improve the two revival bounds. On the box,
the first term in \eqref{tt:angle} is nonpositive. In the factored
feedback difference its first numerator has negative first factor;
the other factor is positive. Its absolute forcing in $u'$ is at most
\[
 \frac2L\,4\,10\,a\,\frac72\,\frac{16}{15}\,\frac32
 =\frac{448a}{L}<\frac{512a}{L}.
\]
Here $|h\sin\theta-a\cos\theta|\leq a$,
$h\cos\theta+a\sin\theta<7/2$, $b/s<4$, and $D\geq1$.
The second-parent feedback term is at most $6|V|u$ because
$b/\sin s<4$, $\sin(su)/s\leq u$, and the second denominator
exceeds $2/3$. Therefore
\[
 u'\leq512a/L+(3/16)u,\qquad
 u(\tau)\leq(512a/L)e^{3/16}<640a/L<1/1024.
\]
The elementary estimate $e^{3/16}<5/4$ follows, for example,
from $e^x\leq(1-x)^{-1}$ for $0\leq x<1$.
For $Z=-V\geq0$, the integrating factor in its equation gives
\[
 Z(\tau)\leq\int_0^\tau
      F\sqrt D\frac{\sin(su)}{\sin s}\,dr
 <11\int_0^\tau u\,dr\leq7040a/L<1/32,
\]
since $3\sqrt{11}(16/15)<11$ and $s/\sin s<16/15$.
Every boundary has now been strictly improved, or has its invariant
sign. The original positive data and integrating factors keep the
positive quantities away from zero. The vector field remains in a
compact subset of its smooth domain; its integral equation and the
local existence theorem therefore extend it through $\tau=1$.
This proves the full prescribed interval, not only a small initial
piece. Once $\eta$ is positive on an interval, the strictly positive
forcing at $u=0$ gives $u>0$ thereafter, and the integral for $Z$
then gives $V<0$. No instantaneous derivative at a flat endpoint is
used to assert this revival.

For the numerator, $E_1,E_2,x,y,c,C>0$ and $D\geq1$ show
\[
 0\leq K_1^2E_1y+bK_2^2DE_2\leq K_2^2D N_3.
\]
Equation \eqref{tt:numerator_derivatives} and its integrating factor
give the first part of \eqref{tt:coefficientvariation}.
For the length, the exact polar transport equation gives
\[
 \left|(\log|\zeta_3|)'\right|
 \leq\epsilon\omega|\sin(\phi_3-\phi_1)|
       +\frac{|V|}{\sin s}|\sin(\phi_3-\phi_2)|.
\]
The determinant coordinates give
\[
 |\sin(\phi_3-\phi_1)|=
 \frac{ag+hb}{\sqrt D\sqrt{g^2+b^2}}\leq a+b/g<7a,
 \qquad
 |\sin(\phi_3-\phi_2)|=\frac b{\sqrt{g^2+b^2}}<6s.
\]
The cosine factors in the exact polar equation have absolute value
at most one; these inequalities thus retain its full contribution.
Consequently
\[
 |(\log|\zeta_3|)'|
 \leq140a/L+(32/5)|V|\leq45196a/L<46000a/L<1/128.
\]
The last numerical inequality uses $a\leq1/512,L\geq16384$.
This proves the second part of \eqref{tt:coefficientvariation}, and
in particular $R(t)<2e^{1/64}<3$.

For the newest linear pair define the exact scaled coefficients
\[
 \mathcal K=\frac{N_3}{\sigma_2^2\sin s\,R},\qquad
 \mathcal B=\frac{\zeta_{3,1}}{\sin s}.
\]
At entry $1/5\leq\mathcal K(t_a)\leq16/15$.
The proved variations give
\[
 \frac16<\frac15e^{-13/128}\leq\mathcal K
 \leq\frac{16}{15}e^{1/64}<\frac98.
\]
For the lower strict inequality use $e^{-13/128}\geq115/128>5/6$;
for the upper use $e^{1/64}\leq64/63$.
Furthermore $(\log\sqrt D)'=ah\epsilon\omega/D\leq22a/L<1/64$.
Because the second phase is vertical at entry,
\[
 \frac{\zeta_{3,1}(t)}{\zeta_{3,1}(t_a)}
 =\frac{\sqrt{D(t_a)}}{\sqrt D}
        \left(\cos\theta+\frac gb\sin\theta\right).
\]
Its lower bound is $(63/64)^2$, and its upper bound is
$1+6u\leq1+3840a/L<9/8$, using $b\geq\sin s$.
Thus
\begin{equation}\label{tt:newestcoefficientbox}
 1/6<\mathcal K<9/8,\qquad 15/32<\mathcal B<9/8.
\end{equation}
Construct the Riccati equation $v'=\mathcal B-\mathcal K v^2$
with its actual entry value. Its field points strictly inward at
$v=1/2$ and $v=3$, since
$15/32-(9/8)/4=3/16>0$ and $9/8-9/6=-3/8<0$.
It exists on the entire interval with $1/2\leq v\leq3$.
Then define
\begin{equation}\label{tt:amplitudeinverse}
 P_3(t)=P_3(t_a)
 \exp\!\left(\int_0^\tau(\mathcal K v-\delta_3R)\,dr\right),
 \qquad \Theta_3=-P_3,\quad
 \Omega_3=-\frac{K_3}{\sigma_2}P_3v.
\end{equation}
Differentiating these expressions gives exactly
$\dot\Theta_3=N_3\Omega_3/(K_3R)-dK_3^2R\Theta_3$
and $\dot\Omega_3=K_3\zeta_{3,1}\Theta_3-dK_3^2R\Omega_3$.
The inverse is thus the unique physical linear pair, by its integral
equation; it has not assumed that its temperature avoids zero.
Finally $\mathcal K v\geq1/12$, $\delta_3R<3/108$, and
$\mathcal K v<27/8<4$ give
$1/18\leq\mathcal K v-\delta_3R\leq4$.
Integrating over the full unit transition proves \eqref{tt:gain}.
\end{proof}

The exact first, second and third diffusion exponents during this
transition are respectively $\delta_1$, $\delta_2\int_0^1D\,d\tau$
and $\delta_3\int_0^1R\,d\tau$. They equal
$dK_j^2\int_{t_a}^{t_1}|\zeta_j(t)|^2dt$ in physical coordinates.
The second phase and amplitude continue evolving inside these
integrals. The end control is exactly $\dot\alpha=\mathcal F_3$,
so the source next steering starts smoothly, with its actual values
$Z_3(t_1)=\zeta_{3,1}(t_1)$ and $|\zeta_3(t_1)|$.
The released periodic profile is not a Laplacian eigenfunction;
realizing these damped trajectories with its compact spatial fields
still requires the explicitly calculated profile and localization
forces. None of the determinant or transition results discards those
forces or asserts an unforced spatial superposition.
